\documentclass[10pt,a4paper]{amsart}
\usepackage[margin=19mm]{geometry}
\usepackage{amsmath,amssymb,mathtools}
\usepackage[scr=rsfs,cal=euler]{mathalfa}
\usepackage{xfrac}
\usepackage[unicode,hidelinks,bookmarks=false]{hyperref}
\newcommand{\ie}{i.e.\ }
\newcommand{\cf}{cf.\ }
\newcommand{\resp}{respectively\ }
\newcommand{\Subsection}{\subsection}
\providecommand{\nobreakdash}{\nobreak\hbox{-}}
\setlength{\emergencystretch}{2em}
\multlinegap0pt
\usepackage{bm}
\usepackage{bbm}
\usepackage{dsfont}
\usepackage{xspace}
\usepackage{cancel}
\usepackage{mathtools}
\usepackage{amsthm}
\makeatletter
\@ifundefined{thm}{\newtheorem{thm}{Theorem}}{}
\makeatother
\title[Composition operators on some semi-Hilbert spaces]{Composition operators on some semi-Hilbert spaces}
\author{Y. Estaremi, M. S. Al Ghafri}
\date{Source: arXiv:2508.05042v1 (CC BY 4.0)}
\begin{document}
\maketitle

\noindent\emph{Source and licence.} Composition operators on some semi-Hilbert spaces. Version arXiv:2508.05042v1 (CC BY 4.0), distributed under the Creative Commons Attribution 4.0 International licence, as recorded in the arXiv licence metadata for this exact version and shown on the abstract page https://arxiv.org/abs/2508.05042v1. Full rights notice: licenses/arxiv-2508.05042-CC-BY-4.0.txt.\par\smallskip

\noindent\textit{Editorial scope.} Counted unit: the Thm whose source label is \texttt{t15} in the article (source0.tex lines 292--308, source label \texttt{t15}), delivered together with the article's own complete original proof of it (source0.tex lines 311--367, one \texttt{proof} environment of 57 lines). No mathematics is written, repaired or re-derived: statement and proof are the article's own text line for line. Required context retained uncounted: none. Every definition, display and label of those lines is retained unchanged. Omitted from this package and not redistributed: 1--291, 309--310, 368--860. Nothing inside a retained excerpt is omitted or altered: every retained body is delivered exactly as the article's lines are written, so no omission receipt of any kind accompanies this package. Citations: the retained excerpts carry no citation command at all, so no bibliography entry of the article is transcribed and no reference is redistributed. Reference adaptation: none was needed and none was made; every reference inside the delivered unit resolves inside it, so no unresolved reference or citation is produced. Printed slips kept literal: no printed word, symbol, hypothesis or number of the counted statement or of its proof was altered. Layout and class: the article's own class is \texttt{amsart} with packages including latexsym, amsmath, amssymb; the wrapper is the lane's pinned \texttt{amsart} editorial wrapper at 10pt on A4, which restates the theorem-like environments the retained text needs and adds the article's own macro definitions that the retained text needs (none), with extra wrapper packages (bm, bbm, dsfont, xspace, cancel, mathtools). The delivered numbering is the unit's own. Assets: no retained span contains a figure environment, a graphics-inclusion command, a PDF-page inclusion, a table or a TikZ picture; no image file is copied into this package, the package declares no asset dependency and no image file is redistributed with it. Licence: version arXiv:2508.05042v1 of this article is distributed under the Creative Commons Attribution 4.0 International licence, as recorded in the arXiv licence metadata for this exact version and shown on the abstract page https://arxiv.org/abs/2508.05042v1. A full rights notice is delivered at licenses/arxiv-2508.05042-CC-BY-4.0.txt. Characters: every retained excerpt is pure ASCII, so the delivered engine, XeLaTeX with its default Latin Modern text font, reports no missing character.
\begin{thm}\label{t15}
Let \( M_u \) be a positive multiplication operator and \( C_{\varphi} \) a bounded composition operator on the Hilbert space \( L^2(\mu) \). Then the following statements hold:

\begin{itemize}
    \item[(a)] The operator \( C_{\varphi} \) is \( M_u \)-normal if and only if
    \[
    \left( \frac{1}{u} \chi_{S(u)} h \right) \circ \varphi \cdot u = \frac{1}{u} \chi_{S(u)} \cdot J,
    \]
    where \( \chi_{S(u)} \) denotes the characteristic function of the support \( S(u) = \{ x \in X : u(x) \neq 0 \} \), and \( J = h_{\varphi}  \cdot E(u) \circ \varphi^{-1} \).

    \item[(b)] The operator \( C_{\varphi} \) is \( M_u \)-quasi-normal if and only if
    \[
    \frac{1}{u} \chi_{S(u)} \cdot J \cdot h = \left( \frac{1}{u} \chi_{S(u)} \cdot J \right) \circ \varphi^{-1} \cdot h,
    \]
    where \( J = h_{\varphi}  \cdot E(u) \circ \varphi^{-1} \).
\end{itemize}
\end{thm}

\begin{proof}
(a) As we know, \( M_u^\dagger = M_{\frac{1}{u} \chi_{S(u)}} \) and \( C_\varphi^* = M_h C_{\varphi^{-1}} E \). Therefore,
\[
C_\varphi^\sharp = M_u^\dagger C_\varphi^* M_u = M_{\frac{1}{u} \chi_{S(u)}} M_h C_{\varphi^{-1}} E M_u.
\]
By definition, \( C_\varphi \) is \( M_u \)-normal if and only if \( C_\varphi^\sharp C_\varphi = C_\varphi C_\varphi^\sharp \). For every \( f \in L^2(\mu) \), we compute:
\[
C_\varphi^\sharp C_\varphi f = M_{\frac{1}{u} \chi_{S(u)}} M_h C_{\varphi^{-1}} E M_u C_\varphi f = \frac{1}{u} \chi_{S(u)} J \cdot f,
\]
\[
C_\varphi C_\varphi^\sharp f = C_\varphi M_{\frac{1}{u} \chi_{S(u)}} M_h C_{\varphi^{-1}} E M_u f = \left( \frac{1}{u} \chi_{S(u)} h \right) \circ \varphi \cdot E(uf).
\]
Hence,
\[
C_\varphi^\sharp C_\varphi f = C_\varphi C_\varphi^\sharp f \quad \text{if and only if} \quad \frac{1}{u} \chi_{S(u)} J \cdot f = \left( \frac{1}{u} \chi_{S(u)} h \right) \circ \varphi \cdot E(uf).
\]
This is equivalent to the integral equality
$$\int_X\frac{1}{u}\chi_{S(u)}J.fd\mu=\int_X(\frac{1}{u}\chi_{S(u)}h)\circ\varphi E(uf)d\mu=\int_X(\frac{1}{u}\chi_{S(u)}h)\circ\varphi ufd\mu,$$
for all \( f \in L^2(\mu) \). Therefore, \( C_\varphi \) is \( M_u \)-normal if and only if
\[
\left( \frac{1}{u} \chi_{S(u)} h \right) \circ \varphi \cdot u = \frac{1}{u} \chi_{S(u)} \cdot J.
\]

\medskip

(b) By definition, \( C_\varphi \) is \( M_u \)-quasi-normal if and only if
\[
C_\varphi^\sharp C_{\varphi^2} = C_\varphi C_\varphi^\sharp C_\varphi.
\]
For every \( f \in L^2(\mu) \), we compute:
\[
C_\varphi^\sharp C_{\varphi^2} f = M_{\frac{1}{u} \chi_{S(u)}} M_h C_{\varphi^{-1}} E M_u C_{\varphi^2} f = \frac{1}{u} \chi_{S(u)} J \cdot f \circ \varphi,
\]
\[
C_\varphi C_\varphi^\sharp C_\varphi f = C_\varphi M_{\frac{1}{u} \chi_{S(u)}} M_h C_{\varphi^{-1}} E M_u C_\varphi f = \left( \frac{1}{u} \chi_{S(u)} \cdot J \right) \circ \varphi \cdot f \circ \varphi.
\]
Thus,
\[
C_\varphi^\sharp C_{\varphi^2} f = C_\varphi C_\varphi^\sharp C_\varphi f
\]
if and only if
\[
\frac{1}{u} \chi_{S(u)} J \cdot f \circ \varphi = \left( \frac{1}{u} \chi_{S(u)} \cdot J \right) \circ \varphi \cdot f \circ \varphi.
\]
Equivalently, we obtain
\begin{align*}
\int_X \left( \frac{1}{u} \chi_{S(u)} J \right) \circ \varphi^{-1} \cdot h \cdot f \, d\mu 
&= \int_X \frac{1}{u} \chi_{S(u)} J \cdot f \circ \varphi \, d\mu \\
&= \int_X \left( \frac{1}{u} \chi_{S(u)} \cdot J \right) \circ \varphi \cdot f \circ \varphi \, d\mu \\
&= \int_X \frac{1}{u} \chi_{S(u)} \cdot J \cdot h \cdot f \, d\mu.
\end{align*}

for all \( f \in L^2(\mu) \). Therefore, \( C_\varphi \) is \( M_u \)-quasi-normal if and only if
\[
\frac{1}{u} \chi_{S(u)} J \cdot h = \left( \frac{1}{u} \chi_{S(u)} \cdot J \right) \circ \varphi^{-1} \cdot h.
\]
\end{proof}

\end{document}
